\section{Threats to Validity}
We discuss the threats to validity of our research following the categorization of Runeson and Höst~\cite{Runeson2009}.
\subsection{Construct Validity}

The construct \emph{security-relevance} is operationalized through expert judgment rather than as an approximation of an objective ground truth. 
Different expert pools might apply a different decision boundary, and the $13.5\%$ of C1 items with at least one disagreeing assessor confirm that the construct admits legitimate boundary cases, as discussed in ~\Cref{sec:discussion-c1}.
We mitigate this by multi-rating every released item, and by reporting Fleiss' $\kappa$ and raw observed agreement rather than a single reliability number.
For C3, \emph{retrieval usefulness} is characterized in two layers: an artifact-level analysis of shortlists (applicability relative to the retrieving item's requirement) and a practitioner $0$--$5$ rating used as calibration of that analysis.

\subsection{Internal Validity}

C2 thresholds and hyperparameters are selected on a single C1 validation split.
The tight spread of F2 across the sentence-encoder variants bounds the risk of threshold-to-split overfitting but does not eliminate it, and probability calibration is not performed beyond threshold choice because C2 is used downstream as a hard filter into C3.
The held-out C1 test set contains only $58$ items, so the reported in-distribution point estimates probably carry wide confidence intervals and, taken alone, cannot rule out overfitting to the two source projects.
We mitigate this by not relying on the in-distribution split for the headline claim. Deployment selection is restricted to the transfer-oriented candidate set, and the zero-shot evaluation on the five Wu et al.\ benchmarks~\cite{wu2022data-quality}, spanning several thousand items from unrelated projects, provides the primary evidence that the learned notion of security relevance generalizes.
The C3 evaluation is preliminary and performed over two industrial product backlogs and two practitioner interviews.
Observations are descriptive and the apparent convergences across the two sessions are consistent with, but do not establish generalizable empirical evidence.
The complete evaluation, involving more practitioners per backlog, multiple backlogs, ideally across organizations, and a study design capable of measuring workflow impact rather than perceived relevance, is deferred to future work.
Item-level familiarity is recorded as descriptive context.
The interviews target the engineer-consulting scenario in which the consumer of a security requirement is not necessarily the author of the corresponding backlog item.
Prior familiarity with a specific backlog item does not, however, disqualify a practitioner's judgment of whether a retrieved security requirement is applicable.

\subsection{External Validity}

C1 is drawn from two Atlassian open-source projects whose item language resembles industrial enterprise backlogs. It is a precision resource by design and does not aim to cover the broader variance of public issue trackers (Wu et al.'s benchmarks fill that role).
The nine C1 annotators share an organizational context, which could bias the operational definition of \texttt{security-relevant}.
We argue in Section~\ref{sec:discussion-c1} that this context is operationally shaped by the same standards other security-regulated organizations follow.

\subsection{Reliability}

Both the C3 requirement extractor and the semantic filter are open-weights LLMs, and both stages introduce run-to-run variability in the resulting shortlist.
We mitigate this by fixing the model identifier, sampling temperature, and prompts, and by reporting averaged behavior.
We do not claim bitwise reproducibility of the C3 output.
Although we do not quantify run-to-run variance formally, repeated executions during development informally suggested that the C3 pipeline was practically consistent: for the same backlog item and source corpus, it reliably returned the same clauses and only small differences in relevance scores.
C1 and the full C2 pipeline are fully reproducible from the released artifacts and are deterministic up to library-level floating-point variation.
The C3 cross-distribution probe on CIS Benchmarks is reproducible end-to-end against public documents. The internal-policy evaluation is reproducible only on the internal corpus.
